% !TeX program = xelatex
\documentclass[12pt,a4paper]{article}
\usepackage{fontspec}
\setmainfont{Times New Roman}
\usepackage[english,russian]{babel}
\usepackage{amsmath,amssymb,mathrsfs,cancel}
\usepackage[margin=22mm]{geometry}
\usepackage{graphicx,float}
\usepackage[unicode,hidelinks]{hyperref}
\DeclareMathOperator{\grad}{grad}
\DeclareMathOperator{\Div}{div}
\newcommand{\Rey}{\mathrm{Re}}
\setlength{\emergencystretch}{3em}
\interfootnotelinepenalty=10000
\newcommand{\topic}[1]{\par\medskip\noindent\textbf{#1}\par\nobreak\smallskip}
\renewcommand{\thefootnote}{\arabic{footnote})}
\begin{document}
\section*{Вопрос 6}
\noindent Гидравлическая теория безнапорной фильтрации. Уравнение Буссинеска. Формула Дюпюи для безнапорной фильтрации.

\medskip
\noindent\textbf{Источник:} Н. Е. Леонтьев. \textit{Основы теории фильтрации}. 2-е изд. Москва: МАКС Пресс, 2017. Глава II, § 9, печатные страницы 54--58.
\par\bigskip
\begin{quote}\sloppy
Строгое математическое описание безнапорных течений\footnote{Как уже говорилось на стр. 43, \textit{безнапорными} называются течения со свободной поверхностью.}, часто встречающихся в гидрологической практике, достаточно сложно из-за неизвестности формы поверхности депрессии. Для приближенного описания течений со свободной границей французским гидравликом и экономистом Жюлем Дюпюи [19] были предложены упрощающие предположения, которые выполняются для течений с \textit{пологой} поверхностью депрессии.


Рассмотрим (для простоты плоское) нестационарное безнапорное течение несжимаемой жидкости в пористом слое, ограниченном снизу горизонтальной непроницаемой поверхностью (рис.~19). Движение жидкости описывается уравнением неразрывности и законом Дарси, которые в показанной на рисунке системе координат имеют вид\footnote{Для сокращения записи мы будем иногда обозначать частные или обыкновенные производные штрихом, указывая переменную, по которой производится дифференцирование, с помощью нижнего индекса, например: $\frac{\partial p}{\partial x}=p'_x$.}
\begin{equation}\begin{gathered}
\frac{\partial u_x}{\partial x}+\frac{\partial u_z}{\partial z}=0,\qquad u_x=-C\frac{\partial}{\partial x}\left(\frac{p}{\rho g}+z\right)=-C\frac{p'_x}{\rho g},\\
u_z=-C\frac{\partial}{\partial z}\left(\frac{p}{\rho g}+z\right)=-C\left(\frac{p'_z}{\rho g}+1\right).
\end{gathered}\tag{9.1}\label{q6:eq-9-1}\end{equation}
Основное предположение гидравлической теории заключается в том, что при малых углах наклона свободной поверхности к горизонту ($h'_x\ll1$) в любом вертикальном сечении пласта $x=\mathrm{const}$ давление в фильтрующейся жидкости распределено по гидростатическому закону (что эквивалентно постоянству напора по сечению). Если эта гипотеза выполняется, то из закона Дарси следует,
что скорость фильтрации в любой точке направлена горизонтально (ее вертикальная компонента равна нулю), а ее величина определяется только формой свободной поверхности и не зависит от высоты:
\[u_x(x,t)=-Ch'_x,\qquad u_z=0.\]
\par\nobreak\smallskip{\small Н. Е. Леонтьев. \textit{Основы теории фильтрации}. 2-е изд., 2017. Глава II, § 9, с. 54--55.}\par
\end{quote}

\begin{quote}\sloppy
С использованием предположений гидравлической теории легко получить уравнение для высоты свободной поверхности $h(x,t)$, записав закон
сохранения массы для части пористой среды между двумя произвольными сечениями $x_1=\mathrm{const}$ и $x_2=\mathrm{const}$ (рис.~19):
\[\begin{aligned}
\frac{d}{dt}\int_{x_1}^{x_2}\rho mh\,dx&=\left.(\rho h\cdot Ch'_x)\right|_{x_1}^{x_2}=\left.\rho C\left(\frac{h^2}{2}\right)'_x\right|_{x_1}^{x_2}\\
&=C\rho\int_{x_1}^{x_2}\left(\frac{h^2}{2}\right)''_{xx}\,dx,
\end{aligned}\]
откуда получается нелинейное \textit{уравнение Буссинеска} (J. Boussinesq, 1904)
\begin{equation}\frac{\partial h}{\partial t}=\frac{C}{2m}\frac{\partial^2(h^2)}{\partial x^2}.\tag{9.3}\label{q6:eq-9-3}\end{equation}
В общем случае зависимости высоты свободной поверхности от \textit{двух} пространственных координат в горизонтальной плоскости уравнение для $h=h(x,y,t)$ выглядит аналогично:
\[\frac{\partial h}{\partial t}=\frac{C}{2m}\triangle(h^2).\]
В частном случае стационарных течений форма поверхности депрессии находится из уравнения Лапласа
\[\triangle(h^2)=0.\]
\par\nobreak\smallskip{\small Н. Е. Леонтьев. \textit{Основы теории фильтрации}. 2-е изд., 2017. Глава II, § 9, с. 55--56.}\par
\end{quote}

\begin{quote}\sloppy
В случае стационарного плоскопараллельного течения уравнение Буссинеска дает для формы свободной поверхности зависимость
\begin{equation}h^2(x)=Ax+B,\qquad A,B=\mathrm{const}\tag{9.4}\label{q6:eq-9-4}\end{equation}

(для течения слева направо, показанного на рис. 19, $A<0$), для которой постоянный объемный расход жидкости (на единицу длины в поперечном направлении)
\[Q=-Ch'(x)\cdot h(x)=-C(h^2(x)/2)'_x=-CA/2.\]
Записав (9.4) для двух сечений $x_1$ и $x_2$, несложно получить выражение расхода через высоты свободной поверхности в этих точках:
\begin{equation}Q=\frac{C}{2}\frac{h^2(x_1)-h^2(x_2)}{x_2-x_1}.\tag{9.5}\label{q6:eq-9-5}\end{equation}
Аналогичным образом для стационарного осесимметричного безнапорного течения получим решение
\[h^2(r)=A\ln\left(\frac{r}{r_0}\right),\quad A,r_0=\mathrm{const}\]
с расходом
\[Q=-Ch'(r)\cdot2\pi rh(r)=-2\pi rC(h^2(r)/2)'_r=-\pi CA\]
(для притока к скважине, показанного на рис.~20, $A>0$ и расход $Q$ отрицательный). Аналогично плоскому случаю, для расхода можно записать выражение
\begin{equation}Q=-\pi C\frac{h^2(r_2)-h^2(r_1)}{\ln(r_2/r_1)}.\tag{9.6}\label{q6:eq-9-6}\end{equation}
Соотношения (9.5) и (9.6) называются \textit{формулами Дюпюи} для пологой безнапорной фильтрации\setcounter{footnote}{0}\footnote{Сравните с формулой Дюпюи (7.5) на стр. 42.}.
\par\nobreak\smallskip{\small Н. Е. Леонтьев. \textit{Основы теории фильтрации}. 2-е изд., 2017. Глава II, § 9, с. 56--57.}\par
\end{quote}

\begin{quote}\sloppy
Формулы Дюпюи являются приближенными, т.к. при их выводе использовалась приближенная теория. Однако, что любопытно, в случае стационарной фильтрации через прямоугольную перемычку (рис.~21) расход, определяемый по гидравлической теории, совпадает с точным значением, причем при любых (даже немалых) наклонах поверхности депрессии.

\par\nobreak\smallskip{\small Н. Е. Леонтьев. \textit{Основы теории фильтрации}. 2-е изд., 2017. Глава II, § 9, с. 57.}\par
\end{quote}

\end{document}
